% =============================================================================
%  Back cover
% =============================================================================
\clearpage
\thispagestyle{empty}
\pagecolor{gsvblue}
\afterpage{\nopagecolor}

\vspace*{\fill}

\begin{center}
\begin{tikzpicture}
\node[text width=12cm,align=center,text=white] {%
    {\fontsize{14}{17}\selectfont\bfseries
     Assessment of Railroad Ballast Degradation}\\[6pt]
    {\fontsize{11}{14}\selectfont\itshape
     under modified drop-weight impact loading conditions\\
     and image-based morphological analysis}\\[20pt]
    {\color{gsvaccent}\rule{6cm}{0.5pt}}\\[14pt]
    {\fontsize{10}{13}\selectfont
     A purpose-built \SI{50}{\kilo\gram} drop-weight apparatus
     was used to subject railway ballast to controlled impact at
     two energy levels (\SI{147}{\joule} and \SI{221}{\joule})
     across nine factorial configurations of pad condition and
     base compliance. Paired pre-and-post-impact 3D scanning
     of 45~tracer particles yielded the first per-particle
     morphological change dataset for Indian railway ballast.
     The Indian Railway Fouling Index $\mathit{FI}_{\mathrm{IN}}$
     spans a 6.4-fold range; polyurethane mats reduce fines
     generation at both energy levels. Mode~A (corner abrasion)
     accounts for the majority of degradation events; Mode~B
     (catastrophic fracture) is concentrated on
     high-energy--rigid-base tests.}\\[20pt]
    {\color{gsvaccent}\rule{3cm}{0.5pt}}\\[14pt]
    {\fontsize{9}{11}\selectfont\itshape
     Gati Shakti Vishwavidyalaya\\
     School of Civil Engineering, Vadodara\\
     April 2026}
};
\end{tikzpicture}
\end{center}

\vspace*{\fill}
